\chapter{Transport System API} \label{transportSystemApi}
CTaps is not a complete implementation of the TAPS standard,
as that would be too large of a scope for a singular thesis.
\todo{Check in on this later, describe how large of a subset}
This chapter will describe which parts of the TAPS API that has
been implemented in CTaps, and what operations that are supported
in CTaps.

\section{Basic TAPS Objects}
The basic TAPS objects:
\texttt{Preonnection}, \texttt{Connection}, \texttt{Local Endpoint}, \texttt{Remote Endpoint}, 
\texttt{Listener} and \texttt{Message} are each represented as a type defined
struct.

\begin{figure}[H]
\begin{minted}{C}
typedef struct Connection {
  TransportProperties transport_properties;
  LocalEndpoint local_endpoint;
  RemoteEndpoint remote_endpoint;
  // TODO - decide on if this has to be a pointer
  ProtocolImplementation protocol;
  uv_handle_t* protocol_uv_handle;
  ConnectionOpenType open_type;
  struct SocketManager* socket_manager;
  // TODO this is shared state and should be locked
  // Queue for pending receive() calls that arrived before the data
  GQueue* received_callbacks;
  GQueue* received_messages;
} Connection;

\end{minted}
\caption{Connection typedef structure}
\end{figure} \todo{Update this when having finalized Connection struct}

Each of these objects has a set of associated functions, in order to encapsulate
the responsibililty of this structure as best as possible. 

\begin{figure}[H]
\begin{minted}{C}
int send_message(Connection* connection, Message* message);
int receive_message(
    Connection* connection,
    ReceiveMessageRequest receive_message_cb
);
void connection_build_from_listener(
    Connection* connection,
    struct Listener* listener,
    RemoteEndpoint* remote_endpoint
);
void connection_close(Connection* connection);
\end{minted}
\caption{Function declarations from connection.h}
\end{figure}

\section{Event Loop}

\section{Asynchrony} \label{asynchrony}
\todo{Can I write some numbers here, how much is I/O bound as a percentage?}
Most modern interactions with networking applications are asynchronous
and the TAPS API reflects this. CTaps is therefore callback based. Any
interactions with the network or other I/O-bound activities are loaded
performed by the central event loop, which is managed by libuv. Whenever
a user wants to initiate a Connection, send a message or receive a message,
they pass a structure containing the relevant function pointers for the given
functionality. \todo{This struct way is not the current implementation, but perhaps the desired one?}
This includes function pointers for when a relevant action succeeds or fails.
If the initiation is successful, then the callback passed to the initiation
function which is associated with a successful initiation is invoked. \todo{Kronglete lol}
When invoking this callback, the new Connection is passed as a parameter,
along with a void pointer consisting of data specified by the user.
This void pointer allows the user to pass any data desired to the callback.

If initiating a preconnection fails because the endpoint could not be
reached, then the function pointer associated with the failure of creating
a connection is invoked. \todo{This is not currently implemented, show code when it has been}
As parameters to this callback, CTaps passes pointer to the relevant (pre?) connection \todo {sort out uncertainity}
and a reason for the failure, letting the user know what caused the error,
so that they can respond accordingly.

Due to the asynchronous nature of the library, users can request to
receive messages at any time, and the library can receive messages
over the network and cache them at any time. When a user requests
to receive a message they pass a callback to the application, if a
message is available, that is, the library has received a message
over the network and has cached it, the callback is invoked immediately,
passing the first cached message as a parameter to the callback.
If no message has been received by the library, then the callback
itself is cached. The callback is then invoked the next time a message
is received over the network. This way, the user does not have to
spin while waiting for messages and is perform other work while
waiting for incoming messages. Each invocation only receives
a single message, and will always receive the first message in
the message queue.

When sending messages the message is not cached, instead
immediately passed to the libuv event loop. The message
may queue here for a while, depending on the number of tasks
in the loop.

\todo{Overview of callbacks here}


\section{Selecting the desired protocol}
Talk about how the user can use selection properties to select the desired protocol 

\section{Control Flow}
Since the primary functionality of the library is handled
by the underlying libuv event loop, the control flow has
to be handed over to the event loop after the primary setup
of an application has been handled. Further control has
to be
handled through the callbacks invoked by the library.

\section{Events}
As mentioned in \cref{asynchrony}, Taps is an asynchronous library.
For CTaps, this means that any asynchronous functions take a struct
containing function pointers as a parameter. In RFC 9622, each
asynchronous routine finishing invokes an Event \cite{trammellAbstractApplicationProgramming2025}.
The notation used in TAPS to denote events is shown in \cref{fig:ctapsEventNotation}.
A full list of events and the objects that spawn them can
be seen in \cref{fig:tapsEvents}. For each Event, CTaps
features an associated callback function.

\begin{figure}[H]
\centering
\begin{BVerbatim}
Object -> Event<>
\end{BVerbatim}
\caption{CTaps Event notation of Object spawning an Event}
\label{fig:ctapsEventNotation}
\end{figure}

This associated callback is set at various times. Mostly,
they are set when initially building the object associated
with the Event. For exapmle, the Ready event is set
when initializing a preconnection. \todo{For listeners as well or?}
In this case, the callback structure is stored in the
corresponding object, so the callback can be invoked when
the given Event occurs.
However, some Events are passed when inititating an asynchronous action,
such as CloneError, where the function pointer is passed as a parameter
to the Connection.clone() function. \todo{Make sure that this is true, will we implement clone?}

In order to allow the user to pass data into the callback, the structure
containing callbacks also contains a void pointer. This void pointer
is then passed as a parameter to the callback. This allows the user
implementation of the callback to dereference the void pointer
and pass any data to the callback, to make it accessible in
the body of the callback itself. This is a standard 
way if passing data to callbacks in C and can be
seen for example in the libuv library, where the
\texttt{uv\_handle\_t} struct contains a void pointer which
can be dereferenced in any callbacks \cite{libuvUv_handle_tDocumentation2025}.

\begin{figure}[H]
    \centering
\begin{minted}{C}
typedef struct ConnectionCallbacks {
  int jalla;
} ConnectionCallbacks;
\end{minted}
\label{fig:connectionCallbackStructure}
\caption{The structure passed when initiating a Preconnection}
\end{figure}

\begin{figure}[H]
\begin{center}
\begin{tabular}{ ||l|c|c| } 
 \hline
 \textbf{Event}      & \textbf{Class}                    & \textbf{Description} \\ 
 \hline \hline
 CloneError          & Connection                        & cell9 \\ 
 \hline
 Closed              & Connection                        & cell9 \\ 
 \hline
 ConnectionError     & Connection                        & cell9 \\ 
 \hline
 ConnectionReceived  & Listener                          & cell9 \\ 
 \hline
 EstablishmentError  & Connection/Preconnection/Listener & cell9 \\ 
 \hline
 Expired             & Connection                        & cell9 \\ 
 \hline
 PathChange          & Connection                        & cell9 \\ 
 \hline
 Ready               & Connection                        & cell6 \\ 
 \hline
 ReceiveError        & Connection                        & cell9 \\ 
 \hline
 ReceivedPartial     & Connection                        & cell9 \\ 
 \hline
 RendezousDone       & Preconnection                     & cell9 \\ 
 \hline
 SendError           & Connection                        & cell9 \\ 
 \hline
 Sent                & Connection                        & cell9 \\ 
 \hline
 SoftError           & Connection                        & cell9 \\ 
 \hline
 Stopped             & Listener                          & cell9 \\ 
 \hline
\end{tabular}
\end{center}
\caption{Table of events and associated objects}
\label{fig:tapsEvents}
\end{figure}
